%%%%%%%%%%%%%%%%%%%%%%%%%%%%%%%%%%%%%%%%%
% Short Sectioned Assignment LaTeX Template Version 1.0 (5/5/12)
% This template has been downloaded from: http://www.LaTeXTemplates.com
% Original author:  Frits Wenneker (http://www.howtotex.com)
% License: CC BY-NC-SA 3.0 (http://creativecommons.org/licenses/by-nc-sa/3.0/)
%%%%%%%%%%%%%%%%%%%%%%%%%%%%%%%%%%%%%%%%%

% \documentclass[paper=a4, fontsize=11pt]{scrartcl} % A4 paper and 11pt font size
\documentclass[11pt]{book}
\usepackage[a4paper, margin=2cm]{geometry}
\usepackage[T1]{fontenc} % Use 8-bit encoding that has 256 glyphs
\usepackage[utf8]{inputenc}
\usepackage{listings} % para insertar código con formato similar al editor
\usepackage[spanish, es-tabla]{babel} % Selecciona el español para palabras introducidas automáticamente, p.ej. "septiembre" en la fecha y especifica que se use la palabra Tabla en vez de Cuadro
\usepackage{url} % ,href} %para incluir URLs e hipervínculos dentro del texto (aunque hay que instalar href)
\usepackage{graphics,graphicx, float} %para incluir imágenes y colocarlas
\usepackage[gen]{eurosym} %para incluir el símbolo del euro
\usepackage{cite} %para incluir citas del archivo <nombre>.bib
\usepackage{enumerate}
\usepackage{hyperref}
\usepackage{graphicx}
\graphicspath{{imágenes/}}
\usepackage{tabularx}
\usepackage{booktabs}
% Permite hacer citacion de bloque
\usepackage{setspace}
\usepackage{etoolbox}

\newcommand{\titulo}{MaquiTracker}
\newcommand{\subtitulo}{Servicio de localización, estadística y mantenimiento preventivo para maquinaria agrícola.}
\newcommand{\tituloingles}{MaquiTracker}
\newcommand{\subtituloingles}{Location, statistics and preventive maintenance service for agricultural machinery.}
\newcommand{\autor}{Daniel Pedrosa Montes}
\newcommand{\tutor}{Juan Julián Merelo Guervos}
\newcommand{\departamento}{Departamento de Ingeniería de Computadores, Automática y Robótica}

\usepackage[table,xcdraw]{xcolor}
\hypersetup{
	colorlinks=true,	% false: boxed links; true: colored links
	linkcolor=black,	% color of internal links
	urlcolor=cyan		% color of external links
}

% Configuración tipografía
\renewcommand{\familydefault}{\sfdefault}
\usepackage[scaled=1]{helvet}
\usepackage[helvet]{sfmath}
\everymath={\sf}

% Header y footer
\usepackage{fancyhdr} % Custom headers and footers
\pagestyle{fancyplain} % Makes all pages in the document conform to the custom headers and footers
\fancyhead[L]{} % Empty left header
\fancyhead[C]{} % Empty center header
\fancyhead[R]{\autor} % My name
\fancyfoot[L]{} % Empty left footer
\fancyfoot[C]{} % Empty center footer
\fancyfoot[R]{\thepage} % Page numbering for right footer
%\renewcommand{\headrulewidth}{0pt} % Remove header underlines
\renewcommand{\footrulewidth}{0pt} % Remove footer underlines
\setlength{\headheight}{13.6pt} % Customize the height of the header

\usepackage{titlesec, blindtext, color}
\definecolor{gray75}{gray}{0.75}
\newcommand{\hsp}{\hspace{20pt}}
\titleformat{\chapter}[hang]{\Huge\bfseries}{\thechapter\hsp\textcolor{gray75}{|}\hsp}{0pt}{\Huge\bfseries}
\setcounter{secnumdepth}{4}
\usepackage[Lenny]{fncychap}

\begin{document}
	% Plantilla portada UGR
	\begin{titlepage}
\newlength{\centeroffset}
\setlength{\centeroffset}{-0.5\oddsidemargin}
\addtolength{\centeroffset}{0.5\evensidemargin}
\thispagestyle{empty}

\noindent\hspace*{\centeroffset}\begin{minipage}{\textwidth}

\centering
\includegraphics[width=0.9\textwidth]{logos/logo_ugr.jpg}\\[1.4cm]

\textsc{ \Large TRABAJO FIN DE GRADO\\[0.2cm]}
\textsc{ GRADO EN INGENIERIA INFORMATICA}\\[1cm]

{\Huge\bfseries \titulo \\}
\noindent\rule[-1ex]{\textwidth}{3pt}\\[3.5ex]
{\large\bfseries \subtitulo }
\end{minipage}

\vspace{2.5cm}
\noindent\hspace*{\centeroffset}
\begin{minipage}{\textwidth}
\centering

\textbf{Autor}\\ \autor \\[2.5ex]
\textbf{Director}\\ {\tutor}\\[2cm]
\includegraphics[width=0.3\textwidth]{logos/etsiit_logo.png}\\[0.1cm]
\textsc{Escuela Técnica Superior de Ingenierías Informática y de Telecomunicación}\\
\textsc{---}\\
Granada, Junio de 2024
\end{minipage}
\end{titlepage}


	% Plantilla prefacio UGR
	\thispagestyle{empty}

\begin{center}
{\large\bfseries \titulo \\ \subtitulo }\\
\end{center}
\begin{center}
\titulo \\
\end{center}

%\vspace{0.7cm}

\vspace{0.5cm}
\noindent\textbf{Palabras clave}: \textit{software libre}
\vspace{0.7cm}

\noindent\textbf{Resumen}\\
	

\cleardoublepage

\begin{center}
	{\large\bfseries \tituloingles \\ \subtituloingles}\\
\end{center}
\begin{center}
	\autor\\
\end{center}
\vspace{0.5cm}
\noindent\textbf{Keywords}: \textit{open source}, \textit{floss}
\vspace{0.7cm}

\noindent\textbf{Abstract}\\


\cleardoublepage

\thispagestyle{empty}

\noindent\rule[-1ex]{\textwidth}{2pt}\\[4.5ex]

D. \textbf{\tutor}, Profesor(a) del \departamento.

\vspace{0.5cm}

\textbf{Informo:}

\vspace{0.5cm}

Que el presente trabajo, titulado \textit{\textbf{\titulo}},
ha sido realizado bajo mi supervisión por \textbf{\autor}, y autorizo la defensa de dicho trabajo ante el tribunal
que corresponda.

\vspace{0.5cm}

Y para que conste, expiden y firman el presente informe en Granada a Junio de 2018.

\vspace{1cm}

\textbf{El/la director(a)/es: }

\vspace{5cm}

\noindent \textbf{\tutor}

\chapter*{Agradecimientos}






	% Índice de contenidos
	\newpage
	\tableofcontents

	% Introducción
	\chapter{Introducción}

Este Trabajo de Fin de Grado consiste en construir una solución software
capaz de emitir, monitorizar y registrar una serie de eventos en un
sistema de control hidráulico embebido. 

Este proyecto es software libre, y está liberado con
la licencia \textbf{GNU Affero General Public License} \cite{agplv3}.

\section{Motivación}

En España el mercado de los cultivos leñosos supone un ??? del PIB
anual, siendo un sector que emplea a más de ??? personas. [cita] Entre estos
predominan el cultivo de olivos, almendros, pistachos, cítricos y
viñedos. [cita] 

La explotación de estas plantaciones se realiza mayoritariamente por personas
físicas y no por empresas. [cita] La profesionalización de los propietarios
es creciente, pero baja. [cita] A su vez ocurre que no existe una renovación
de los jefes de explotación, con un porcentaje de gestores menores de 45 años
inferior al 40\%. La formación reglada específica de los jefes en el caso de
las explotaciones agrarias es inferior al 25\% en todas las comarcas españolas,
llegando a menos del ??? en aquellas regiones con mayor
explotación de este tipo de cultivos. [cita] Encontramos la mayor superficie
de cultivos leñosos en la mitad sur de la península, destacando Andalucía,
Castilla la Mancha y Extremadura. [cita INE https://storymaps.arcgis.com/stories/6fa8c26ecbde4fd99a0446e874c64898]

Desgraciadamente, el bajo nivel de automatización y de inversión por parte
de los propietarios de las explotaciones agrarias hace que la rentabilidad
de estos haya disminuido en los últimos años. [cita]
El sector de la
mecanización de este tipo de cultivos ha crecido en los últimos años de una
forma lenta y con algunas pausas, dado el alto coste de la maquinaria, la
falta de emprendimiento en el sector y el hecho de que la mayoría de las
explotaciones son de pequeño tamaño. [cita] 

Dos terceras partes del
coste de este tipo de cultivos van destinados a la recolección y a la poda.
[cita https://centroliva.com/mecanizacion-de-la-recoleccion-del-olivar-desde-la-investigacion-hasta-la-industria/]
Entre las soluciones mecanizadas
para la recolección encontramos principalmente los vibradores y los paraguas,
mientras que para la poda encontramos las tijeras eléctricas, las sierras y
las trituradoras. [cita]

[ Crear o extraer gráficos del INE que ayuden a entender los datos aquí expuestos ]

Son muchos los jefes de explotación que deciden externalizar la recolección
de sus cultivos a empresas especializadas. [cita?] Ocurre un problema en este
caso, y es que el jefe de explotación deja de tener la certeza de que la
recolección se ha realizado de forma correcta, y de que la cantidad de producto
recolectado es la que se espera. Este problema sucede también en otros
contextos, como el de grandes fincas profesionalizadas.


\section{Objetivo}

Este proyecto busca desarrollar una solución software que permita
registrar, exportar y visualizar información sobre la recolección de cultivos
leñosos con paraguas vibradores según el criterio de nuestro cliente,
con la restricción de que el registro y el exportado pueda ejecutarse en
una plataforma hardware ya distribuida en el mercado.

Entre otros objetivos secundarios se encuentran:

\begin{enumerate}
    \item Automatización de los procesos de desarrollo y despliegue de la solución, como es
    la generación de documentación, la ejecución de pruebas, la generación de librerías
    y el despliegue en la plataforma hardware.
    \item Escribir código que cumpla los estándares de la industria para maquinaria pesada,
    asegurando que no puedan ocurrir situaciones inesperadas en tiempo de ejecución que puedan
    comprometer la integridad del código de control, que se ejecutará sobre la misma plataforma
    hardware.
    \item Generar una documentación adecuada para que cualquier desarrollador pueda entender,
    mantener y ampliar este proyecto en el futuro.
\end{enumerate}

\section{Estructura del documento}

[ Explica qué secciones tiene el documento y qué desarrollan ]

[ Explica el por qué de la elección ]

	
	\input{secciones/02_metodología}
	% Descripción del problema y hasta donde se llega
	\chapter{Problema a resolver}

\textbf{Buscamos describir de una forma clara y concisa el problema y sus
circunstancias}. Consideramos que el problema se describe de forma clara y
concisa cuando se enuncia con el mínimo nivel de complejidad posible%
%
\footnote{Aunque
no por ello debe ser insuficiente. Un mito extendido acerca de las metodologías ágiles
es que no necesitamos tener una idea clara del problema antes de empezar a desarrollar
una solución. Nada mas lejos de la verdad, lo que se busca con los métodos ágiles en la
fase de análisis es poner a prueba empíricamente que nuestro modelo del problema es fiel
a la realidad, y que el problema que estamos resolviendo es el correcto antes de proponer
una solución arquitectónica. Así se concibe en las publicaciones que iniciaron el hoy
degenerado movimiento. [Cita Martin]
}%
, el lenguaje que se utiliza es simple pero adecuado
en su contexto y no tiene sesgos hacia ninguna solución. 

Formulamos el problema haciendo uso de las fases de empatizar y definir del método
\textit{Design Thinking} para después, tanto en un análisis previo como a medida
que vamos desarrollando el proyecto, extraer historias de usuario, historias de
desarrollador y \textit{user journeys} de donde obtenemos los objetivos deseados
junto a sus respectivas importancias.
Finalmente sintetizamos esta información en forma de requisitos funcionales y
requisitos no funcionales.

\section{Método}

Es inmensamente doloroso el poco tiempo que vivimos, pero es que es más doloroso dedicar
nuestro tiempo a tareas que no benefician a aquellos entes a los cuales queremos aportar
valor.
\textbf{A lo largo de este Trabajo de Fin de Grado exponemos diversos métodos para recibir
realimentación y para conseguirla en el menor tiempo posible}, de forma que minimicemos
el tiempo invertido en tareas sin valor para nuestro propósito.

Este primer método lo ideamos una vez el cliente nos comunica brevemente su necesidad.
\textbf{Su principal objetivo es el de recuperar y modelar el problema a resolver}, ya sea en
el análisis previo o como complementario del análisis continuo, pero también busca
otros objetivos como el de transmitir confianza al cliente o el de ayudarnos como empresa a
decidir si esta será una batalla en la que nos gustaría e interesaría participar.

\textbf{Como primer paso tenemos que el cliente comunique su necesidad.} Sólo por el hecho de que esta
persona física o jurídica está depositando confianza en nosotros merece el más
adecuado respeto, que se demuestra dejándole el espacio que necesita para exponer
sus ideas. De esta descripción extraemos información verbal y
no verbal que nos ayuda a comprender las circunstancias en que ocurre el potencial
problema a resolver, y en caso de no ser un cliente de confianza o desconocer las
circunstancias en que se desenvuelve el problema, si como empresa nos interesa
ser un proveedor de soluciones para él y su sector.

De una forma atenta, pero crítica, escuchamos y anotamos la información que el
cliente quiere transmitirnos. Si existe algún término que enuncia y desconocemos
o alguna circunstancia que nos describe y que no llegamos a comprender o bien le
interrumpimos o bien anotamos para preguntar más adelante.

\textbf{El segundo paso es identificar la raíz del problema.} Por experiencia sabemos
que las personas tienden a enunciar los problemas en términos de la solución%
%
\footnote{[ Code Complete 2 ]
}%
. Muchas veces aquello que nos solicitan responde a una necesidad subyacente que no
se identifica o, si identificada, no se comunica.

Buscamos dialogando con el cliente y por medio de nuestra propia investigación exponer
el problema eliminando las restricciones innecesarias, buscando aquellas que son
necesarias y, a poder ser, sin omitir aspectos que pudiesen afectar a la arquitectura
de la solución.

Por ejemplo, proponemos borradores rápidos de arquitectura los cuales
nuestro cliente puede analizar de forma crítica, de forma independiente
o por medio del método dialéctico. El interesado desmiente aquellos aspectos
pertinentes y saca a la luz nuevas necesidades y restricciones que no se habían
comunicado.

[ Tercer paso: formalización de las conversaciones por medio de las herramientas.
Explica que lo hacemos por ejemplificar, pero que el uso de estas herramientas
dependerá en nuestros conocimientos del sector, en el tamaño y toma del contacto
del sector de nuestro equipo, en la complejidad del problema y en el nivel de intuición
que podamos tener en este ]

\subsection{Herramientas}

[ Enuncia problema que resuelven las historias de usuario ]

[ Describe cómo resuelve el problema la historia de usuario ]

[ Describe la esencia de una historia de usuario ]

[ Bridge desde la historia de usuario hacia las historias de desarrollador ]

[ Igual que con las historias de usuario pero para las user stories y
los requisitos funcionales y no funcionales ]

\section{Análisis previo}

[ Explica los objetivos del análisis previo ]

\subsection{Enunciado del problema haciendo uso de la técnica del \textit{Design Thinking}}

[ Explica en términos generales Design Thinking. ]

[ Explica por qué la técnica del Design Thinking puede ayudarnos a este propósito ]

[ Explica cómo adaptamos la técnica a nuestras necesidades. ]

\subsubsection{Conocimientos sobre los usuarios y el problema general (Comprender)}

Existen empresas y trabajadores los cuales, a cambio de dinero, trabajan la
tierra de un tercero con su maquinaria.

\textbf{Quien contrata el servicio} quiere reportes/resúmenes acerca de la actividad
de la máquina dado un intervalo de tiempo. Quiere conocer, a partir de una fecha y hora de inicio y otra de final:

\begin{enumerate}
   \item Tiempo medio y mediano entre vibraciones.
   \item Tiempo medio y mediano vibrando.
   \item Posición de los árboles vibrados.
   \item Tiempo de trabajo efectivo.
\end{enumerate}

\textbf{Tanto quien contrata el servicio como quien lo ejecuta} quieren tener una seguridad alta
en que los reportes/resúmenes son verdaderos y que no han sido modificados.

\textbf{El software tiene que poder ejecutarse en una plataforma hardware ya distribuida en el mercado.}
La maquinaria se controla mediante un sistema de control dedicado con las siguientes características:

\begin{itemize}
   \item Renesas RA6M1.
   \begin{itemize}
      \item Arm® Cortex®-M4.
      \item 8kB Data FLASH.
      \item 128kB SRAM.
      \item USB2.0 FS.
   \end{itemize}
   \item Conexión mediante UART a módulo de comunicaciones.
   \begin{itemize}
      \item BMD-380. (Sin firmware).
      \item ESP32-C3. (Firmware desarrollado: comunicación emulando UART a través de
      BLE).
      \item Digi XBee Pro. (Firmware de fábrica: UART a través de ZigBee.)
   \end{itemize}
   \item FRAM Fujitsu 4kiB. (Posibilidad de otros tamaños de almacenamiento). Driver
   desarrollado.
   \item Posibilidad de incluir otra FRAM Fujitsu o FLASH que sea compatible con la
   interfaz del integrado.
   \item Smart GNSS antenna module UBLOX CAM-M8Q. Driver desarrollado.
   \item Entradas 4-20mA. Driver desarrollado.
   \item Entradas digitales. Driver desarrollado.
   \item Entradas CAN 2.0. Driver desarrollado.
   \item Alimentación a 12V.
\end{itemize}

\subsubsection{Empatía con los usuarios mirándoles de cerca (Observar)}

Cualquier solución que se proponga debería de incluir interfaces
descriptivas y amigables para personas sin conocimientos de
electricidad o electrónica.

\subsubsection{Usuario típico (Definir el punto de vista)}

Empresa o trabajador que se dedica a tratar los árboles de campos a un tercero.

Suelen ser hombres que están acostumbrados al trabajo físico. Operan maquinaria
agrícola. Para comunicarse telemáticamente suelen utilizar Whatsapp y e-mail.

[ **Pendiente - Investigar hasta qué punto están dispuestos a voluntariamente
echar a andar la aplicación cada vez que quieran registrar un trabajo** ]

\subsection{Historias de usuario}

\subsection{Historias de desarrollador}

\subsubsection{HD??: feedback estimaciones}
Para mejorar la calidad de futuras estimaciones quiero registrar datos acerca del desarrollo
para poder:

\begin{itemize}
   \item Comparar las estimaciones de tiempo del segundo capítulo con los datos reales.
   \item Comparar las estimaciones de tiempo del segundo capítulo con los datos reales.
\end{itemize}

\subsection{\textit{User Journeys}}

\subsection{Síntesis}

[ Requisitos Funcionales ]

[ Requisitos No funcionales ]

\section{Análisis continuo}

\subsection{Historias de usuario}

\subsection{Historias de desarrollador}

\subsection{\textit{User Journeys}}

\subsection{Síntesis}

[ Requisitos Funcionales ]

[ Requisitos No funcionales ]


	% Arquitectura del sistema
	% 1. Análisis del problema
	% 2. Posibles arquitecturas
	% 3. Arquitectura elegida junto al cliente
	\chapter{Arquitectura del sistema}

\section{Posibles soluciones}

\subsection{Generar todas las ideas posibles (Idear)}

\subsubsection{Histórico de trabajo (almacenamiento en control electrónico)}

Aunque podría ser más lento de recuperar y de desarrollar potencialmente
es más conveniente para el cliente.

\begin{enumerate}
   \item Almacenamiento de los datos.
   \begin{itemize}
       \item Para almacenar se puede utilizar \href{https://github.com/armink/FlashDB}{FlashDB}. 
       Único proyecto vivo profesional que encuentro apto para plataformas baremetal, listo para producción y de código abierto. Está mantenido activamente por una empresa. Permite TSDB y KVDB. Si vemos que no es suficiente al ser de código abierto podemos hacer una librería que dependa en el sistema o un fork con las mejoras en el software y un pull request de vuelta.
   \end{itemize}

   \item Generación de los datos de salida (con garantías).
   \begin{itemize}
       \item Se puede generar un PDF en la tarjeta y firmarlo.
       \item Se puede generar un bloque de texto en la tarjeta y firmarlo.
       \item Se puede generar una estructura binaria y firmarla.
   \end{itemize}

   \item Transmisión de los datos a un \textit{smartphone}, ordenador o similar.
   \begin{itemize}
       \item Mediante UART-BLE. \href{https://stackoverflow.com/a/22919464}{5kB/s~10kB/s}.
       \begin{itemize}
           \item Módulo piggyback ESP-C3 o BDM-380.
           \item Requiere que al otro lado se ejecute un software nuestro.
       \end{itemize}
       \item Mediante WiFi Direct.
       \begin{itemize}
           \item Diseñando un módulo piggyback con algún módulo compatible con este protocolo.
           \item No es compatible con iPhone, al menos no de una forma fácil.
           \begin{itemize}
               \item Requiere una aplicación.
           \end{itemize}
           \item No sabemos el porcentaje de móviles Android que lo implementan.
       \end{itemize}
       \item Mediante WiFi.
       \begin{itemize}
           \item Diseñando un módulo piggyback con algún módulo compatible con este protocolo.
           \begin{itemize}
               \item Requiere implementar un servidor SFTP, HTTP, FTP o similar.
           \end{itemize}
           \item Requiere que al otro lado se ejecute un software nuestro.
       \end{itemize}
       \item Mediante Bluetooth File Transfer profile.
       \begin{itemize}
           \item Diseñando un módulo piggyback con algún módulo compatible con este protocolo.
           \item No es compatible con iPhone.
       \end{itemize}
       \item Mediante Bluetooth Personal Area Network profile.
       \begin{itemize}
           \item Diseñando un módulo piggyback con algún módulo compatible con este protocolo.
           \item Es compatible con el iPhone sin necesidad de hacerse parte del programa MFi.
           \item Requiere que al otro lado se ejecute un software nuestro.
       \end{itemize}
       \item Mediante USB.
       \begin{itemize}
           \item Requiere utilizar algún stack que nos permita utilizar el protocolo USB.
           \item Requiere utilizar algún stack que nos permita interactuar con un sistema de ficheros FAT o similar.
           \item Requiere adaptar físicamente la placa y la caja, garantizando que este agujero no permita que entre polvo al control.
       \end{itemize}
   \end{itemize}
\end{enumerate}

\subsubsection{Histórico de trabajo (almacenamiento en *smartphone*)}

Aunque requeriría que el cliente estuviese atento de iniciar y finalizar
una sesión de trabajo desde su teléfono, potencialmente puede ser más
rápido y barato de desarrollar.

Parte de la premisa de que lo que se quiere extraer de aquí es la información
de un día de trabajo y no el histórico de una máquina.

\begin{enumerate}
    \item Transmisión de datos entre control y *smartphone*. (con garantías)
    \begin{itemize}
        \item BLE Peripheral Publisher.
        \item ¿Utilizando el procedimiento \href{https://lpccs-docs.renesas.com/Tutorial-DA145x-BLE-Security/access_and_signing.html}{Connection Data Signing}.?
        \item Necesitamos una aplicación de móvil que se suscriba a una característica.
        \item En iOS el servicio parece funcionar en primer plano y en segundo plano
              te suscribes a una característica en concreto.
        \item En Android va a ser más complicado de testar, pero en principio
              rece más o menos estable.
        \item Necesitamos tener un buffer en el control para el caso de que se
    desconecte el dispositivo, así como un aviso por pantalla de que se
    ha perdido la conexión.
    \end{itemize}
    \item Almacenamiento de datos. (con garantías)
    \begin{itemize}
        \item Firestore/Realtime Database. Implementado Android e iOS,
        opcionalmente puede sincronizarse con una base de datos en
        Google Cloud. El framework de Firebase ofrece otros servicios.
        De pago.
        \item SQLite. Gestor de bases de datos SQL. Amplia adopción y de código
        abierto.
        \item MongoDB. Gestor de bases de datos NoSQL. Amplia adopción y de código
        abierto.
    \end{itemize}
\end{enumerate}

\subsubsection{Histórico de trabajo (visualización de los datos)}

[ REPLANTEAR COMPLETAMENTE Y AÑADIR MÁS OPCIONES MÁS ESPECÍFICAS ]

1. Comprobación de la validez de los datos y visualización.
    - Mediante un servicio web que permita visualizar los datos en la
    propia página web.
        - Requeriría un mantenimiento durante muchos años por nuestra parte.
    - Mediante un servicio web que permita transformar los datos a un PDF.
        - Requeriría un mantenimiento durante muchos años por nuestra parte.
    - Mediante un software que permita transformar los datos a un PDF.
        - Requeriría distintos ports a distintas plataformas.
    - Mediante un software que permita visualizar y consultar los datos.
        - Requeriría distintos ports a distintas plataformas.
    - Mezcla de alguna de las anteriores.

	
	\chapter{Implementación}

La implementación del software se ha dividido en hitos. Estos, han sido definidos en GitHub
y cada uno de ellos contiene un grupo de \textit{issues} que se corresponden con las distintas
mejoras que se han ido incorporando al software a lo largo de su desarrollo.\\



	\chapter{Comparación de los datos previos, durante y posteriores al desarrollo}

[ Compara el problema inicial planteado en el tercer capítulo con el problema final. ]

[ Compara la solución propuesta en el cuarto capítulo con la solución final. ]

[ Compara las estimaciones de tiempo del segundo capítulo con los datos reales. ]

[ Compara las estimaciones de tiempo del segundo capítulo con los datos reales. ]



	% Conclusion
	\chapter{Conclusión}



	\newpage
	\bibliography{bibliografia}
	\bibliographystyle{plain}

	% Índice de imágenes y tablas
	\listoffigures

	% Si hay suficientes se incluirá dicho índice
	\listoftables
	\newpage
\end{document}
